% !TEX root = main.tex


Strings are ubiquitous in modern software systems, especially with the
advent of programming languages like JavaScript, PHP, and Python. Despite this, 
string manipulation is well-known to be error-prone and can easily lead to 
security vulnerabilities including HTML injection (e.g.\ see \cite{Kern14,LB16}).
Applications to analysis of security vulnerabilities caused by string
manipulation have been one of the main catalysts for the extensive research into
SMT over strings spanning across the last twenty years
\cite{LB16,DBLP:conf/cav/AbdullaACHRRS15,ostrich,Berkeley-JavaScript,HAMPI,BTV09,string-survey}.
One of the success stories of string solvers includes their usage at AWS
for analysis of Role-Based Access Control (RBAC) policies (e.g.\ see
\cite{neha,Cook18}).
Since 2020 SMT-LIB Unicode String theory \cite{SMTLIB} has been formalized and supported by
many existing string solvers, including Z3 \cite{z3}, Z3-alpha \cite{z3alpha}, Z3-Noodler \cite{noodler,noodler-int,noodler-len, noodler-tool}, Z3str3RE \cite{z3str3re},
cvc5 \cite{cvc5}, Z3str4 \cite{z3str4}, Trau \cite{trau0,trau-tool,trau2}, and our own solver OSTRICH \cite{ostrich}.

In this paper, we present OSTRICH2: the latest evolution of the string solver 
OSTRICH \cite{ostrich}. The goal behind the development of OSTRICH in 2018 was 
to address the lack of support for complex string functions like replaceall, 
transducers, reverse, etc. by string solvers. In particular, these are known to
be crucial for capturing real-world string manipulations
\cite{Berkeley-JavaScript,framework,LB16}. In a nutshell, this initial version
of OSTRICH employed the so-called \emph{backward regular propagation} for the
so-called straight-line fragments, and leveraged the use of \emph{Nielsen's
transformation}, when the formula is not straight-line. Over the years, 
OSTRICH has undergone some substantial changes, which have radically enhanced
the solver, resulting (among others) in the superior performance of OSTRICH in
competitions (e.g. \#1 for QF\_S track in SMT-Competition 2023). This paper aims
to, for the first time, introduce potential future users to OSTRICH in its
latest form. To this end, rather than being exhaustive, we aim for 
accessibility by taking the reader through illustrative examples (among others).
In addition to unravelling the design of OSTRICH2, we also report our latest
experimentation of the solver on SMT-COMP'24 and SMT-COMP'25 benchmarks, showing
its competitiveness, most notably on unsatisfiable instances.

The following is a list of main features  of OSTRICH2:
\begin{description}
    \item[(F1)] SMT-LIB Unicode String theory inputs.
    \item[(F2)] Native support of features of ECMAScript regular expressions (e.g. 
        lookaround assertions, capture groups, references).
    \item[(F3)] Native support of replaceall and, more generally, complex string
        transductions using a more comprehensive \emph{regular constraint
        propagation} strategy, which includes not only backward, but also 
        \emph{forward} regular propagations (among others).
    \item[(F4)] An extensive portfolio of solving strategies that includes
%The original OSTRICH string‐solver framework—built atop Princess with automata‐based reasoning—was introduced in \cite{ostrich}. In subsequent work we enriched it with prioritized transducers and full ECMA-2020 regex parsing (look-around assertions, capture groups) \cite{popl22}. OSTRICH2 is a complete rewrite that unifies three core solving engines—Regular Constraint Propagation (RCP), 
        a cost-enriched string solver (CE-Str) \cite{atva2020}, a list-based ADT 
        solver (ADT-Str), %—within a portfolio architecture, 
        and preprocessings (length/character-count abstraction).
        %and new inference rules.
\end{description}

\paragraph*{Organization} Section~\ref{sec:grammar} reviews the SMT-LIB 
Unicode Strings standard and our grammar (a fully formal description is given in Appendix~\ref{app:grammar}); Section~\ref{sec:architecture} reviews the architecture of OSTRICH2; Section~\ref{sec:string-theory-algorithms} details OSTRICH2’s key algorithms; Section~\ref{sec:completeness} discusses the completeness of the algorithms; Section~\ref{sec:extensibility} shows that OSTRICH2 is extensible with user-defined string functions; Section~\ref{sec:experiments} reports evaluation results; and Section~\ref{sec:conclusion} concludes with directions for future work.
